% Part: computability
% Chapter: machines-computations
% Section: halting-problem

\documentclass[../../../include/open-logic-section]{subfiles}

\begin{document}

\olfileid{tur}{und}{hal}
\olsection{رُکّن دا مسئلہ}

\begin{explain}
فرض کرو اسی ٹیورنگ مشیناں دے بیانواں دی اک گنتی وار
فہرست مقرر کر لئی اے۔ ایس طرح ہر ٹیورنگ مشین نوں اک
\emph{اشاریہ} ملدا اے: مشین دے بیانواں دی فہرست
$M_1$، $M_2$، $M_3$، \dots{} وچ اوہدے بیان دی جگہ۔

سانوں پتا اے کہ کجھ فنکشناں دا حساب ٹیورنگ مشین نال
نہیں ہو سکدا: مشیناں دے بیانواں دا مجموعہ، تے ایس لئی
ٹیورنگ مشیناں دا مجموعہ، !!{enumerable} اے، پر
$\Nat$ توں $\Nat$ تک سارے فنکشناں دا مجموعہ
ایہو جیہا نہیں۔ پر اسی ناقابلِ حساب فنکشناں دیاں خاص
مثالاں وی لبھ سکدے آں۔ اک مثال رُکّن دا فنکشن اے۔
\end{explain}

\begin{defn}[رُکّن دا فنکشن]
 \emph{رُکّن دا فنکشن}~$h$ ایویں تعریف ہوندا اے:
\[
h(e,n) =
\begin{cases}
  \text{0} & \text{جے مشین~$M_e$ انپٹ $n$ اُتے نہ رُکے} \\
  \text{1} & \text{جے مشین~$M_e$ انپٹ $n$ اُتے رُکے}
\end{cases}
\]
\end{defn}

\begin{defn}[رُکّن دا مسئلہ]
\emph{رُکّن دا مسئلہ} ایہ طے کرن دا مسئلہ اے کہ کسے وی
$e$، $n$ لئی ٹیورنگ مشین~$M_e$ انپٹ وچ~$n$ لکیراں
لے کے رُکدی اے یا نہیں۔
\end{defn}

\begin{explain}
اسی~$h$ دا ٹیورنگ-ناقابلِ حساب ہونا اک متعلقہ فنکشن
~$s$ دا ٹیورنگ-ناقابلِ حساب ہونا دکھا کے ثابت کراں گے۔
ایہ ثبوت ایس گل اُتے ٹکیا اے کہ جیہڑا وی کم ٹیورنگ
مشین حساب کر سکدی اے، اوہ پابند ٹیورنگ مشین وی کر
سکدی اے (\olref[mac][dis]{sec})؛ تے ایس گل اُتے
کہ دو ٹیورنگ مشیناں جوڑ کے اک مشین بنائی جا سکدی اے
(\olref[mac][cmb]{sec})۔
\end{explain}

\begin{defn} فنکشن~$s$ ایویں تعریف ہوندا اے:
\[
s(e) =
\begin{cases}
  \text{0} & \text{جے مشین~$M_e$ انپٹ $e$ اُتے نہ رُکے} \\
  \text{1} & \text{جے مشین~$M_e$ انپٹ $e$ اُتے رُکے}
\end{cases}
\]
\end{defn}

\begin{lem}
فنکشن~$s$ ٹیورنگ مشین نال حساب نہیں ہو سکدا۔
\end{lem}

\begin{proof}
تضاد لئی فرض کرو کہ فنکشن~$s$ ٹیورنگ مشین نال حساب
ہو سکدا اے۔ فیر اک ٹیورنگ مشین~$S$ ہووے گی جیڑی
~$s$ دا حساب کرے۔ عمومیّت گھٹائے بغیر اسی فرض کر
سکدے آں کہ جدوں $S$ رُکدی اے، اوہدے سر تھلّے
پٹی دا پہلا خانہ ہوندا اے (یعنی اسی اوہنوں پابند روپ وچ
لے سکدے آں)۔ ایہ مشین اک دوجی مشین~$J$ نال
“جوڑی” جا سکدی اے۔ جے اوہ انپٹ~$0$ نال شروع
ہووے، یعنی پٹی دے سِرے دے نشان دے سجّے پاسے
پہلے خانے اُتے ابتدائی حالت وچ~$\TMblank$ پڑھے،
تاں اوہ رُک جاندی اے؛ دوجی صورت وچ سجّے پاسے
چلدی رہندی اے تے کدی نہیں رُکدی۔ $S \concat J$،
یعنی $S$ نوں~$J$ نال جوڑ کے بنی مشین، خود ٹیورنگ
مشین اے؛ ایس لئی ایہ $M_e$ اے، کسے نہ کسے~$e$ لئی
(یعنی فہرست وچ کسے نہ کسے تھاں آندی اے)۔ ہن $M_e$
نوں~$e$ عدد $\TMstroke$ لکیراں دے انپٹ اُتے
شروع کرو۔ دو صورتاں نیں: یا تاں $M_e$ رُکدی اے
یا نہیں رُکدی۔
\begin{enumerate}
\item فرض کرو $M_e$ انپٹ وچ $e$ عدد $\TMstroke$
  لکیراں لے کے رُکدی اے۔ فیر $s(e)
  = 1$ اے۔ ایس لئی $S$، انپٹ~$e$ اُتے شروع
  ہو کے، پٹی اُتے اکّو $\TMstroke$ آؤٹ پٹ دے
  کے رُکدی اے۔ فیر $J$ پٹی اُتے اک $\TMstroke$
  دے نال شروع ہوندی اے۔ ایس صورت وچ $J$ نہیں
  رُکدی۔ پر $M_e$، $S
  \concat J$ مشین اے؛ ایس لئی اوہنوں عین اوہو
  کرنا چاہیدا اے جیڑا $S$ توں بعد $J$ کردیاں نیں
  (یعنی ایتھے سجّے پاسے چلدے رہنا تے کدی نہ رُکنا)۔
  ایس لئی $M_e$ انپٹ وچ $e$ عدد $\TMstroke$
  لکیراں لے کے رُک نہیں سکدی۔

\item ہن فرض کرو $M_e$ انپٹ وچ $e$ عدد
  $\TMstroke$ لکیراں لے کے نہیں رُکدی۔ فیر
  $s(e) = 0$ اے، تے $S$، انپٹ~$e$ اُتے
  شروع ہو کے، خالی پٹی دے نال رُکدی اے۔ $J$
  خالی پٹی اُتے شروع ہو کے فوراً رُکدی اے۔ فیر
  $M_e$ وی اوہو کردی اے جیڑا $S$ توں بعد~$J$
  کردیاں نیں؛ ایس لئی $M_e$ انپٹ وچ $e$ عدد
  $\TMstroke$ لکیراں لے کے ضرور رُکدی اے۔
\end{enumerate}
دونوں صورتاں وچ ساڈے فرض دے خلاف نتیجہ نکلیا۔
ایس توں پتا لگدا اے کہ ایہو جیہی ٹیورنگ مشین~$S$
نہیں ہو سکدی: $s$~ٹیورنگ مشین نال حساب نہیں ہو سکدا۔
\end{proof}

\begin{thm}[رُکّن دے مسئلے دا ناقابلِ حل ہونا]
\ollabel{thm:halting-problem} رُکّن دا مسئلہ حل نہیں ہو
سکدا؛ یعنی فنکشن~$h$ ٹیورنگ مشین نال حساب نہیں
ہو سکدا۔
\end{thm}

\begin{proof}
فرض کرو $h$ ٹیورنگ مشین نال حساب ہو سکدا ہووے،
مثلاً اک مشین~$H$ نال۔ اسی $H$ ورت کے~$s$ دا
حساب کرن والی مشین بنا سکدے آں: پہلے انپٹ دی
اک نقل بنا لو (تے دوہاں نوں اک~$\TMblank$ نشان
نال وکھ رکھو)۔ فیر شروع ول واپس آ کے~$H$ چلاؤ۔
پہلا کم کرن والی مشین صاف طور اُتے بن سکدی اے
(\cref{tur:mac:dis:prob:copier})؛ جے $H$ موجود
ہوندی، تاں اسی اوہنوں ایہو جیہی نقل کرن والی مشین نال
“جوڑ” کے نویں مشین بنا لیندے جیڑی طے کردی کہ
$M_e$ انپٹ~$e$ اُتے رُکدی اے یا نہیں، یعنی~$s$
دا حساب کردی۔ پر اسی پہلے ہی دکھا چکّے آں کہ
ایہو جیہی کوئی مشین نہیں ہو سکدی۔ ایس لئی $h$~وی
ٹیورنگ مشین نال حساب نہیں ہو سکدا۔
\end{proof}

\begin{prob}
تِن-رُکّن (3-Halt) دا مسئلہ ایہ اے کہ کسے وی چنی
ہوئی ٹیورنگ مشین لئی ایہ طے کرن دا فیصلہ جاتی
طریقہ دِتا جاوے کہ اوہ باقی خالی پٹی اُتے تِن
$\TMstroke$ لکیراں دے انپٹ نال رُکدی اے یا نہیں۔
ثابت کرو کہ تِن-رُکّن دا مسئلہ حل نہیں ہو سکدا۔
\end{prob}

\begin{prob}
دکھاؤ کہ جے مشین/انپٹ جوڑیاں $M_e$ تے~$n$
لئی، جتھے $e \neq n$، رُکّن دا مسئلہ حل ہو سکدا
ہووے، تاں $e = n$ والی صورتاں لئی وی حل ہو
سکدا اے۔
\end{prob}

\begin{prob}
اسی ثابت کیتا کہ رُکّن دا مسئلہ حل نہیں ہو سکدا جدوں
انپٹ اک عدد~$e$ ہووے، جیڑا ساریاں ٹیورنگ
مشیناں دی فہرست راہیں اک مشین~$M_e$ دی نشاندہی
کردا اے۔ پر جے اسی \olref[tur][und][enu]{sec} وچ دتے
مشیناں دے بیان ای سدھے
انپٹ بنن دی اجازت دیئیے تاں کیہ ہووے گا؟ کی کوئی
ٹیورنگ مشین مشین دے اشاریے دی تھاں اوہدا بیان
انپٹ وچ لے کے رُکّن دا مسئلہ حل کر سکدی اے؟
دسو کیوں یا کیوں نہیں۔
\end{prob}

\begin{prob} دکھاؤ کہ ایویں تعریف شدہ \emph{جزوی}
فنکشن~$s'$ ٹیورنگ مشین نال حساب \emph{ہو سکدا} اے:
  \[
  s'(e) =
  \begin{cases}
    \text{1} & \text{جے مشین~$M_e$ انپٹ $e$ اُتے رُکے}\\
    \text{غیر تعریف شدہ} & \text{جے مشین~$M_e$ انپٹ $e$ اُتے نہ رُکے}
  \end{cases}
  \]
\end{prob}

\end{document}
